
\mychapter{Detailed Hardware \& Software Implementation}{Detailed Hardware \& Software Implementation}{}
\label{chap:Detailed-Hardware-and-Software-Implementation}

% Proposed Page Budget: 9 Pages (Pages 19 - 27)
% Core Objective: Detail the physical construction, component integration, register configurations, and code architectures.

This chapter details the transition from theoretical modeling to physical realization. It justifies the silicon selection, outlines the two-layer printed circuit board (PCB) design and routing strategies, and breaks down the embedded C++ architecture required to drive the analog front-end (AFE) and serve the real-time web interface.

% Focus: Evaluate and justify the selection of the ESP32-S3 and AD5941 based on math constraints from Chapter 3.

% TODO: Detail how the ESP32-S3 meets the RAM budget for on-board DSP arrays and WebSocket streaming.
% TODO: Detail why the AD5941 was selected (e.g., dedicated potentiostat topology, high precision).
% Discuss GPIO availability, SPI bus speeds, and hardware interrupt capabilities linking the two ICs.


\mysection{Choosing a microcontroller}{Choosing a Microcontroller}
\label{sec:choosing_mcu}

In the previous section the requirements of the features that need to be available on the microcontroller were explored. The AD5941 interface requires a fast, dedicated hardware SPI bus and a few GPIO pins for interrupts/resets. The web server requires an integrated Wi-Fi radio, a robust TCP/IP stack, and significant RAM/flash to store and serve HTML/JavaScript files. The onboard display requires either I2C (for simple OLEDs) or a high-speed SPI/parallel interface for colour TFT screens. The edge-\ac{ML} engine requires relatively high clock speeds, hardware floating-point support, and enough memory to store model weights (likely using TensorFlow Lite for Microcontrollers).

\mysubsection{Choosing the microcontroller family}{Choosing the Microcontroller Family}
\label{sec:mcu_family}

The range of microcontroller options is vast; there are three main architecture options available: STM32, Nordic (nRF52), and ESP, which are compared in Table~\ref{tab:mcu_family}.

\begin{table}[H]
    \caption{Comparison of candidate microcontroller families.}
    \footnotesize
    \centering
    \begin{tabular}{p{3cm}p{5.5cm}p{5.5cm}}
    \toprule
        \textbf{Platform} & \textbf{Advantage} & \textbf{Disadvantage} \\
    \midrule
        STM32 & Golden standard for microcontrollers & Additional hardware needed for Wi-Fi \\
        Nordic (nRF52) & Dominates the BLE space & Additional hardware needed for Wi-Fi \\
        ESP & Native Wi-Fi, Bluetooth & --- \\
    \bottomrule
    \end{tabular}
    \label{tab:mcu_family}
\end{table}

Since ESP is the only family of microprocessors with native Wi-Fi, Bluetooth, and sufficient processing power, a decision needs to be made regarding which variant of these chips to select. Table~\ref{tab:esp_variants} compares the relevant ESP32 variants.

\begin{table}[H]
    \caption{Comparison of candidate ESP32 variants.}
    \footnotesize
    \centering
    \begin{tabular}{p{3cm}p{3.5cm}p{3.5cm}p{3.5cm}}
    \toprule
        \textbf{Variant} & \textbf{Core Architecture} & \textbf{AI/ML Hardware Acceleration} & \textbf{Best Fit For} \\
    \midrule
        ESP32 (Classic) & Dual-core Xtensa LX6 & No & General purpose, older designs \\
        ESP32-C3 & Single-core RISC-V & No & Ultra-low power, but too weak for ML and Web \\
        ESP32-S3 & Dual-core Xtensa LX7 & Yes (Vector Instructions) & Heavy compute, displays, ML \\
    \bottomrule
    \end{tabular}
    \label{tab:esp_variants}
\end{table}




\section{Physical Circuitry \& Electrode Interface Design}
% Focus: Schematic design, PCB routing, and physical connections.

% TODO: Present schematic connections between the ESP32-S3 and AD5941.
% TODO: Discuss KiCad routing decisions, specifically ground plane stitching, analog/digital plane splitting, and decoupling capacitor placement to minimize noise.
% TODO: Detail the physical connection to the screen-printed electrodes and active shielding for the high-impedance Reference Electrode line.

\begin{figure}[H]
    \centering
    % \includegraphics[width=0.9\textwidth]{figures/pcb_layout_2layer.png}
    \caption{Custom two-layer PCB layout highlighting ground plane stitching and analog/digital signal separation.}
    \label{fig:pcb_layout}
\end{figure}

\section{Analog Front-End (AFE) Register Configuration}
% Focus: High-resolution configuration of the AD5941 SoC.

% TODO: Detail the low-level register programming.
% Explain configuration of the High-Speed and Low-Power DACs, the AFE sequencer, and the programmable gain amplifier (PGA).
% Explain adjusting the internal hardware filters (Sinc3 and Sinc2) to match the voltammetric sweep rates.

\begin{lstlisting}[caption={Example AD5941 AFE sequencer initialization for a voltammetric sweep.}, label={lst:ad5941_init}]
// TODO: Insert your actual AD5941 sequencer or register initialization code here
void AD5941_Init(void) {
    // 1. Wake up AFE
    AD5941_WakeUp();
    // 2. Configure High-Speed DAC
    // 3. Set PGA gain and Sinc filters
}
\end{lstlisting}

\section{ESP32-S3 Firmware Architecture \& Real-Time Loop}
% Focus: Multi-threaded embedded C++ software using FreeRTOS.

% TODO: Explain task division in FreeRTOS.
% Core 0: High-speed SPI interrupts, DMA buffer management, DSP filtering (Moving Average / Savitzky-Golay).
% Core 1: WebSocket server, Wi-Fi stack management, user inputs.

\begin{lstlisting}[caption={FreeRTOS task pinning to separate cores on the ESP32-S3.}, label={lst:freertos_tasks}]
// TODO: Insert your FreeRTOS xTaskCreatePinnedToCore logic here
xTaskCreatePinnedToCore(
    DSP_SPI_Task,       // Task function
    "DSP_Task",         // Task name
    8192,               // Stack size
    NULL,               // Parameters
    configMAX_PRIORITIES - 1, // High priority
    &Task1,             // Task handle
    0                   // Pinned to Core 0
);

xTaskCreatePinnedToCore(
    WebSocket_UI_Task,
    "UI_Task",
    8192,
    NULL,
    1,                  // Lower priority
    &Task2,
    1                   // Pinned to Core 1
);
\end{lstlisting}

\section{Wireless Web GUI \& Interface Realization}
% Focus: The user-facing software interface served from the ESP32-S3.

% TODO: Detail how the HTML5/CSS3/JS web interface is stored and served locally via SPIFFS flash memory.
% TODO: Explain the JavaScript WebSocket link implementation used to receive, parse, and plot the voltammetric sweeps in real-time.

\begin{figure}[H]
    \centering
    % \includegraphics[width=0.9\textwidth]{figures/gui_dashboard.png}
    \caption{Real-time web interface served from SPIFFS, plotting DPV scans via WebSockets.}
    \label{fig:web_gui}
\end{figure}